\documentclass{amsart}
\usepackage{amsmath,amssymb,amsthm,hyperref}

\newtheorem{theorem}{Theorem}
\newtheorem{lemma}{Lemma}
\newtheorem{proposition}{Proposition}
\newtheorem{corollary}{Corollary}
\theoremstyle{definition}
\newtheorem{definition}{Definition}
\newtheorem{remark}{Remark}

\newcommand{\Ord}{\mathrm{Ord}}
\newcommand{\interp}[1]{[\![#1]\!]}
\newcommand{\To}{\to_\beta}
\newcommand{\FV}{\mathrm{FV}}

\begin{document}

\title{An elementary ordinal assignment witnessing strong normalization of STLC}
\maketitle

\section{Statement}

Let $\Lambda_\to$ be the set of simply typed $\lambda$-terms over a set of base
types, with $\beta$-reduction $\To$. We construct a map
\[
  \# : \Lambda_\to \longrightarrow \Ord
\]
such that
\[
  M \To N \;\Longrightarrow\; \#M > \#N \qquad (M,N\in\Lambda_\to),
\]
where the construction of $\#$ is \emph{elementary}: it is defined by structural
recursion on types and terms and never invokes strong normalization (SN). Since
$\Ord$ is well-founded, the existence of $\#$ yields at once a direct proof of SN:
an infinite reduction $M_0\To M_1\To\cdots$ would give an infinite strictly
descending sequence $\#M_0 > \#M_1 > \cdots$ of ordinals, which is impossible.

The map $\#$ is in fact \emph{natural-number} valued (we land in $\omega\subseteq\Ord$),
which is more than enough; for pure STLC the longest reduction sequence is finite, so
ordinals below $\omega$ suffice. We phrase the codomain as $\Ord$ to match the problem
statement and to make the well-foundedness step explicit.

\section{Conventions on the calculus}

\emph{Simple types} are generated by $\sigma ::= o \mid (\sigma\to\tau)$, where $o$
ranges over a fixed set of base types (the argument treats all base types identically,
so we may assume a single base type $o$). \emph{Terms} are typed as usual; we write
$M:\sigma$ when $M$ has type $\sigma$, work up to $\alpha$-equivalence, write $\FV(M)$
for the free variables, and write $M[x:=N]$ for capture-avoiding substitution. The
reduction relation $\To$ is the \emph{full} compatible (congruence) closure of
\[
  (\lambda x.M)\,N \;\To\; M[x:=N];
\]
a redex may be contracted anywhere inside a term. Nothing below presupposes that $\To$
is normalizing.

\section{The interpretation domains}

For each type $\sigma$ we define a set $\mathcal M_\sigma$ of \emph{functionals}, two
relations $<_\sigma$ (strict) and $\le_\sigma$ (non-strict) on $\mathcal M_\sigma$, and
a \emph{weight} $\|\cdot\|_\sigma:\mathcal M_\sigma\to\mathbb N_{\ge 1}$, all by
recursion on $\sigma$. Write $\mathbb P=\{1,2,3,\dots\}$ with the usual order.

\begin{definition}[Domains, orders, weight]\label{def:dom}
\leavevmode
\begin{itemize}
\item \textbf{Base type.}
  $\mathcal M_o=\mathbb P$, with $m<_o n :\Leftrightarrow m<n$,
  $m\le_o n:\Leftrightarrow m\le n$, and $\|n\|_o:=n$.
\item \textbf{Arrow type.}
  $\mathcal M_{\sigma\to\tau}$ is the set of all functions
  $f:\mathcal M_\sigma\to\mathcal M_\tau$ that are \emph{both strictly and non-strictly
  monotone}, i.e.\ that satisfy, for all $a,a'\in\mathcal M_\sigma$,
  \[
    \text{(strict)}\quad a<_\sigma a' \Rightarrow f(a)<_\tau f(a'),
    \qquad
    \text{(non-strict)}\quad a\le_\sigma a' \Rightarrow f(a)\le_\tau f(a').
  \]
  These functions are ordered everywhere-strict and everywhere-non-strict
  \emph{pointwise}:
  \[
    f<_{\sigma\to\tau}g :\Leftrightarrow
    \forall a\in\mathcal M_\sigma,\ f(a)<_\tau g(a),
    \qquad
    f\le_{\sigma\to\tau}g :\Leftrightarrow
    \forall a\in\mathcal M_\sigma,\ f(a)\le_\tau g(a),
  \]
  with weight read off at a canonical least argument:
  \[
    \|f\|_{\sigma\to\tau}:=\|f(\mathbf 0_\sigma)\|_\tau .
  \]
\end{itemize}
The \emph{canonical element} $\mathbf 0_\sigma$ and the \emph{lifts}
$\mathrm{lift}_\sigma:\mathbb P\to\mathcal M_\sigma$ are defined by
\[
  \mathrm{lift}_o(k):=k,\qquad
  \mathrm{lift}_{\sigma\to\tau}(k):=\bigl(a\mapsto\mathrm{lift}_\tau(k+\|a\|_\sigma)\bigr),
  \qquad \mathbf 0_\sigma:=\mathrm{lift}_\sigma(1).
\]
\end{definition}

\begin{remark}[Why membership bundles \emph{both} monotonicities]\label{rem:bundle}
The decisive structural choice is that \emph{membership in $\mathcal M_{\sigma\to\tau}$
already requires non-strict monotonicity in the argument}, not merely strict
monotonicity. This is genuinely stronger: at an arrow domain a strictly monotone
functional need \emph{not} be non-strictly monotone. For instance, with domain
$\mathcal M_{o\to o}$ the map $g\mapsto 10\,g(\mathbf 0)-g(\mathbf 0')$ (for two fixed
arguments $\mathbf 0<\mathbf 0'$, suitably normalised into $\mathbb P$) is strictly
monotone with respect to the everywhere-strict order yet fails the everywhere-non-strict
order. Bundling both conditions into membership has two payoffs. First, every element of
every $\mathcal M_\tau$ is automatically non-strictly monotone, so \emph{any} value --- in
particular the bound argument $a$ fed to an abstraction --- may be inserted into an
environment with no side condition; this removes a real gap that arises if one only
requires strict monotonicity and then needs non-strict monotonicity of a bound variable
used in function position. Second, the order $\le_\sigma$ must still be the
\emph{pointwise} relation, not ``$<_\sigma$ or $=$'' (Remark~\ref{rem:leq}); the two
genuinely differ at arrow type, and the pointwise version is the right inductive
invariant.
\end{remark}

\begin{remark}[Why $\le_\sigma$ must be pointwise]\label{rem:leq}
It is essential that $\le_\sigma$ be defined \emph{pointwise} (recursively), and
\emph{not} as the relation ``$<_\sigma$ or $=$''. At an arrow type these two notions
genuinely differ: if $f,g\in\mathcal M_{\sigma\to\tau}$ satisfy $f(a)\le_\tau g(a)$ for
all $a$, yet $f(a_0)=g(a_0)$ for some $a_0$ while $f(a_1)<_\tau g(a_1)$ for some other
$a_1$, then $f\le_{\sigma\to\tau}g$ holds while neither $f<_{\sigma\to\tau}g$ nor $f=g$
holds. Such situations arise in the monotonicity argument (Lemma~\ref{lem:official}); the
weaker pointwise $\le_\sigma$ is the right inductive invariant, whereas ``$<$ or $=$''
would not be preserved by the induction.
\end{remark}

\begin{lemma}[Basic properties]\label{lem:basic}
For every type $\sigma$:
\begin{enumerate}
\item[(a)] For each $k\in\mathbb P$, $\mathrm{lift}_\sigma(k)\in\mathcal M_\sigma$, and
  $\mathrm{lift}_\sigma$ is monotone in its integer argument in both senses:
  $j<j'\Rightarrow\mathrm{lift}_\sigma(j)<_\sigma\mathrm{lift}_\sigma(j')$ and
  $j\le j'\Rightarrow\mathrm{lift}_\sigma(j)\le_\sigma\mathrm{lift}_\sigma(j')$.
  In particular $\mathbf 0_\sigma\in\mathcal M_\sigma$.
\item[(b)] The weight is order-preserving in both senses:
  $v<_\sigma v'\Rightarrow\|v\|_\sigma<\|v'\|_\sigma$ and
  $v\le_\sigma v'\Rightarrow\|v\|_\sigma\le\|v'\|_\sigma$.
\item[(c)] $\le_\sigma$ is a partial order; $<_\sigma$ is a strict partial order; and they
  interleave: $v<_\sigma v'\Rightarrow v\le_\sigma v'$; if $v\le_\sigma v'$ and
  $v'<_\sigma v''$ then $v<_\sigma v''$; if $v<_\sigma v'$ and $v'\le_\sigma v''$ then
  $v<_\sigma v''$.
\item[(d)] $<_\sigma$ is well-founded.
\end{enumerate}
\end{lemma}

\begin{proof}
We prove (a), (b), (c) by simultaneous induction on $\sigma$; (d) follows from (b).

\emph{(c).} At base type these are the standard facts about $(\mathbb P,<,\le)$. At arrow
type every assertion is the universally quantified pointwise statement, which holds by the
IH for $\tau$ applied at each argument $a$. (For antisymmetry: $f\le g$ and $g\le f$ mean
$f(a)\le g(a)\le f(a)$, so $f(a)=g(a)$ for all $a$, i.e.\ $f=g$.)

\emph{(b).} Base: $\|n\|_o=n$. Arrow: if $f<_{\sigma\to\tau}g$ then
$f(\mathbf 0_\sigma)<_\tau g(\mathbf 0_\sigma)$, so by IH(b),
$\|f\|=\|f(\mathbf 0_\sigma)\|_\tau<\|g(\mathbf 0_\sigma)\|_\tau=\|g\|$; the non-strict case
is identical with $\le$. (Here $\mathbf 0_\sigma\in\mathcal M_\sigma$ by IH(a).)

\emph{(a).} Base: $\mathrm{lift}_o(k)=k$, and $j<j'\Rightarrow j<j'$, $j\le j'\Rightarrow
j\le j'$. Arrow $\sigma\to\tau$: first we check
$\mathrm{lift}_{\sigma\to\tau}(k)=(a\mapsto\mathrm{lift}_\tau(k+\|a\|_\sigma))\in
\mathcal M_{\sigma\to\tau}$, i.e.\ it is \emph{both} strictly and non-strictly monotone in
$a$.
\emph{Strict:} $a<_\sigma a'\Rightarrow\|a\|_\sigma<\|a'\|_\sigma$ (IH(b)) so
$k+\|a\|_\sigma<k+\|a'\|_\sigma$, hence $\mathrm{lift}_\tau(k+\|a\|_\sigma)<_\tau
\mathrm{lift}_\tau(k+\|a'\|_\sigma)$ by IH(a) for $\tau$ (its integer-monotonicity clause).
\emph{Non-strict:} $a\le_\sigma a'\Rightarrow\|a\|_\sigma\le\|a'\|_\sigma$ (IH(b)) so
$k+\|a\|_\sigma\le k+\|a'\|_\sigma$, hence $\mathrm{lift}_\tau(k+\|a\|_\sigma)\le_\tau
\mathrm{lift}_\tau(k+\|a'\|_\sigma)$ by IH(a) for $\tau$. Thus
$\mathrm{lift}_{\sigma\to\tau}(k)\in\mathcal M_{\sigma\to\tau}$.
Its integer-monotonicity: for $j<j'$ and every $a$, $j+\|a\|_\sigma<j'+\|a\|_\sigma$, so
$\mathrm{lift}_\tau(j+\|a\|_\sigma)<_\tau\mathrm{lift}_\tau(j'+\|a\|_\sigma)$ by IH(a) for
$\tau$; this holds for all $a$, hence $\mathrm{lift}_{\sigma\to\tau}(j)<_{\sigma\to\tau}
\mathrm{lift}_{\sigma\to\tau}(j')$. The non-strict integer-monotonicity is the same with
$\le$.

\emph{(d).} By (b), $\|\cdot\|_\sigma$ maps $(\mathcal M_\sigma,<_\sigma)$
strictly-order-preservingly into the well-founded $(\mathbb P,<)$, so an infinite
$<_\sigma$-descending chain would yield an infinite $<$-descending chain in $\mathbb P$;
hence $<_\sigma$ is well-founded.
\end{proof}

\begin{remark}
Lemma~\ref{lem:basic}(d) is the crucial point flagged in the literature: the general
pointwise order on a function space need not be well-founded. We sidestep this by carrying
the explicit weight $\|\cdot\|_\sigma$, which embeds each $(\mathcal M_\sigma,<_\sigma)$
strictly-order-preservingly into the well-founded $(\mathbb P,<)$. All comparisons we
ultimately make descend to comparisons of natural numbers (Theorem~\ref{thm:main}), with no
reference to reduction or SN.
\end{remark}

\begin{definition}[Weight shift]\label{def:add}
For $k\in\mathbb N$ define $\mathrm{add}_\sigma(\cdot,k):\mathcal M_\sigma\to\mathcal M_\sigma$ by
\[
  \mathrm{add}_o(n,k):=n+k,\qquad
  \mathrm{add}_{\sigma\to\tau}(f,k):=\bigl(a\mapsto\mathrm{add}_\tau(f(a),k)\bigr).
\]
\end{definition}

\begin{lemma}[Properties of $\mathrm{add}$]\label{lem:add}
For every type $\sigma$, all $v,v'\in\mathcal M_\sigma$, $k,k'\in\mathbb N$:
\begin{enumerate}
\item[(a)] $\mathrm{add}_\sigma(v,k)\in\mathcal M_\sigma$;
\item[(b)] $\|\mathrm{add}_\sigma(v,k)\|_\sigma=\|v\|_\sigma+k$;
\item[(c)] $v<_\sigma v'\Rightarrow\mathrm{add}_\sigma(v,k)<_\sigma\mathrm{add}_\sigma(v',k)$,
  and $v\le_\sigma v'\Rightarrow\mathrm{add}_\sigma(v,k)\le_\sigma\mathrm{add}_\sigma(v',k)$;
\item[(d)] $k<k'\Rightarrow\mathrm{add}_\sigma(v,k)<_\sigma\mathrm{add}_\sigma(v,k')$;
  and $\mathrm{add}_\sigma(v,0)=v$.
\end{enumerate}
\end{lemma}

\begin{proof}
Induction on $\sigma$. Base: (a) trivial; (b) $\|n+k\|=n+k$; (c) $n<n'\Rightarrow n+k<n'+k$
and $n\le n'\Rightarrow n+k\le n'+k$; (d) $k<k'\Rightarrow n+k<n+k'$, $n+0=n$.
Arrow $\sigma\to\tau$: (a) we must check $\mathrm{add}(f,k)$ is \emph{both} strictly and
non-strictly monotone. \emph{Strict:}
$a<_\sigma a'\Rightarrow f(a)<_\tau f(a')\Rightarrow_{\text{IH(c)}}
\mathrm{add}_\tau(f(a),k)<_\tau\mathrm{add}_\tau(f(a'),k)$. \emph{Non-strict:}
$a\le_\sigma a'\Rightarrow f(a)\le_\tau f(a')\Rightarrow_{\text{IH(c)}}
\mathrm{add}_\tau(f(a),k)\le_\tau\mathrm{add}_\tau(f(a'),k)$. (Both use that $f\in\mathcal
M_{\sigma\to\tau}$ is strictly \emph{and} non-strictly monotone.) Hence
$\mathrm{add}(f,k)\in\mathcal M_{\sigma\to\tau}$.
(b) $\|\mathrm{add}(f,k)\|=\|\mathrm{add}_\tau(f(\mathbf 0_\sigma),k)\|_\tau=
\|f(\mathbf 0_\sigma)\|_\tau+k=\|f\|+k$ by IH(b). (c),(d) are the pointwise statements via
IH(c),(d); e.g.\ for (c) non-strict, $f\le g$ means $f(a)\le g(a)$ for all $a$, so by IH(c)
$\mathrm{add}_\tau(f(a),k)\le\mathrm{add}_\tau(g(a),k)$ for all $a$, i.e.\
$\mathrm{add}(f,k)\le\mathrm{add}(g,k)$.
\end{proof}

\section{Interpretation of terms}

An \emph{environment} $\rho$ assigns to each variable $x:\tau$ a value
$\rho(x)\in\mathcal M_\tau$; $\rho[x:=a]$ denotes update. \emph{No further condition on
$\rho$ is needed}: every element of every $\mathcal M_\tau$ is, by
Definition~\ref{def:dom}, automatically both strictly and non-strictly monotone.

\begin{definition}[Term interpretation]\label{def:interp}
For $M:\sigma$ and an environment $\rho$ on $\FV(M)$, define $\interp{M}_\rho\in\mathcal M_\sigma$
by recursion on $M$:
\begin{align*}
  \interp{x}_\rho &:= \rho(x),\\
  \interp{M\,N}_\rho &:= \interp{M}_\rho\bigl(\interp{N}_\rho\bigr),\\
  \interp{\lambda x^\sigma.M}_\rho &:=
     \Bigl(a\in\mathcal M_\sigma\ \mapsto\
       \mathrm{add}_\tau\!\bigl(\interp{M}_{\rho[x:=a]},\ 1+\|a\|_\sigma\bigr)\Bigr),
     \quad M:\tau .
\end{align*}
\end{definition}

This is plain structural recursion, manifestly total, using no normalization hypothesis.
The decisive design choice is at abstraction: the result is shifted by $1+\|a\|_\sigma$, an
amount that \emph{strictly increases with the argument $a$}. This makes abstractions
strictly monotone \emph{even when $x$ does not occur in $M$}, and furnishes the strict drop
at a redex.

\begin{remark}[Why $1+\|a\|$ and not a fixed ``$+1$'']\label{rem:why}
A naive choice $\interp{\lambda x.M}_\rho=(a\mapsto\interp{M}_{\rho[x:=a]})$, or with a
constant ``$+1$'', \emph{fails}: if $x\notin\FV(M)$ then $\interp{M}_{\rho[x:=a]}$ is
independent of $a$, so the function is constant, hence not strictly monotone, so it does not
lie in $\mathcal M_{\sigma\to\tau}$. This breaks Lemma~\ref{lem:cong}(ii). Shifting by
$1+\|a\|_\sigma$ makes the dependence on $a$ strict, repairing the gap while still leaving a
strict ``$+1$'' for the redex drop.
\end{remark}

\section{The monotonicity lemma}

This is the technical heart. We prove well-definedness and both environment-monotonicities
by one structural induction. Because membership in $\mathcal M_{\sigma\to\tau}$ already
includes non-strict monotonicity in the argument (Definition~\ref{def:dom},
Remark~\ref{rem:bundle}), the previously delicate clause ``the interpreted functional is
non-strictly monotone'' is now simply \emph{part of well-definedness} (M0) and requires no
separate hypothesis or auxiliary term. In particular environments may carry arbitrary
values from the $\mathcal M_\tau$, with no monotonicity side condition to discharge.

\begin{lemma}[Monotonicity of interpretations]\label{lem:official}
For every term $M:\sigma$ and all environments $\rho,\rho'$ on $\FV(M)$:
\begin{enumerate}
\item[\textup{(M0)}] $\interp{M}_\rho\in\mathcal M_\sigma$ \emph{(}so, if
  $\sigma=\tau'\to\sigma''$, $\interp{M}_\rho$ is both strictly and non-strictly monotone in
  its argument\emph{)};
\item[\textup{(M$\le$)}] if $\rho(y)\le_{\tau_y}\rho'(y)$ for all $y\in\FV(M)$, then
  $\interp{M}_\rho\le_\sigma\interp{M}_{\rho'}$;
\item[\textup{(M$<$)}] if moreover $\rho(x)<_{\tau_x}\rho'(x)$ for some $x\in\FV(M)$ and
  $\rho(y)=\rho'(y)$ for $y\ne x$, then $\interp{M}_\rho<_\sigma\interp{M}_{\rho'}$.
\end{enumerate}
\end{lemma}

\begin{proof}
By induction on $M$. We use Lemma~\ref{lem:basic}(c) (order facts, including
$<\Rightarrow\le$ and the mixed-transitivity rules) and Lemma~\ref{lem:add}(c) (that
$\mathrm{add}(\cdot,k)$ preserves $\le$ and $<$). Throughout, ``$\interp{T}_\theta$ is
strictly/non-strictly monotone in its argument'' for a function-typed subterm $T$ is just
the statement that $\interp{T}_\theta\in\mathcal M$, supplied by IH (M0) for $T$.

\medskip\noindent\emph{Variable $M=y$.}
(M0): $\rho(y)\in\mathcal M_{\tau_y}$ (so it is strictly and non-strictly monotone if
$\tau_y$ is an arrow --- this is exactly membership). (M$\le$): $\rho(y)\le\rho'(y)$.
(M$<$): $x=y$ gives $\rho(y)<\rho'(y)$.

\medskip\noindent\emph{Abstraction $M=\lambda z^{\tau'}.P$, $P:\sigma'$,
$\sigma=\tau'\to\sigma'$.} Rename $z$ fresh. For any $a\in\mathcal M_{\tau'}$ the updated
environment $\rho[z:=a]$ is again an environment (no side condition on $a$), so the IHs for
$P$ apply at $\rho[z:=a]$ for arbitrary $a$.
\begin{itemize}
\item (M0). We show $\interp{M}_\rho$ is both strictly and non-strictly monotone.
  \emph{Strict.} Fix $a<_{\tau'}a'$. By IH \textup{(M$\le$)} for $P$ (environments differ
  only at $z$, where $a\le a'$ by $<\Rightarrow\le$),
  $\interp{P}_{\rho[z:=a]}\le_{\sigma'}\interp{P}_{\rho[z:=a']}$. By Lemma~\ref{lem:basic}(b),
  $\|a\|_{\tau'}<\|a'\|_{\tau'}$. Hence by Lemma~\ref{lem:add}(c),(d) and mixed transitivity,
  \[
    \interp{M}_\rho(a)=\mathrm{add}(\interp{P}_{\rho[z:=a]},1{+}\|a\|)
    \le\mathrm{add}(\interp{P}_{\rho[z:=a']},1{+}\|a\|)
    <\mathrm{add}(\interp{P}_{\rho[z:=a']},1{+}\|a'\|)=\interp{M}_\rho(a').
  \]
  \emph{Non-strict.} Fix $a\le_{\tau'}a'$. By IH \textup{(M$\le$)} for $P$,
  $\interp{P}_{\rho[z:=a]}\le_{\sigma'}\interp{P}_{\rho[z:=a']}$; and
  $\|a\|_{\tau'}\le\|a'\|_{\tau'}$ (Lemma~\ref{lem:basic}(b)), so $1{+}\|a\|\le1{+}\|a'\|$.
  Then by Lemma~\ref{lem:add}(c) (value side, $\le$) and (c)+(d) (shift side, $\le$),
  \[
    \interp{M}_\rho(a)=\mathrm{add}(\interp{P}_{\rho[z:=a]},1{+}\|a\|)
    \le_{\sigma'}\mathrm{add}(\interp{P}_{\rho[z:=a']},1{+}\|a'\|)=\interp{M}_\rho(a').
  \]
  Hence $\interp{M}_\rho\in\mathcal M_{\tau'\to\sigma'}$.
\item (M$\le$)/(M$<$). For each $a$, apply IH \textup{(M$\le$)}/\textup{(M$<$)} for $P$ at
  $\rho[z:=a],\rho'[z:=a]$ (these differ exactly where $\rho,\rho'$ do, since the value at
  $z$ is the common $a$), then $\mathrm{add}(\cdot,1{+}\|a\|)$ (Lemma~\ref{lem:add}(c)).
  Holding for all $a$, this gives $\interp{M}_\rho\le_{\tau'\to\sigma'}\interp{M}_{\rho'}$,
  resp.\ $<$, by the pointwise definitions (Definition~\ref{def:dom}).
\end{itemize}

\medskip\noindent\emph{Application $M=PQ$, $P:\tau'\to\sigma$, $Q:\tau'$.}
Let $p=\interp{P}_\rho,p'=\interp{P}_{\rho'},q=\interp{Q}_\rho,q'=\interp{Q}_{\rho'}$.

(M0): by IH \textup{(M0)} for $P,Q$, $p,p'\in\mathcal M_{\tau'\to\sigma}$ and
$q,q'\in\mathcal M_{\tau'}$, so $p(q)\in\mathcal M_\sigma$ (apply the strictly-and-non-strictly
monotone $p$ to $q$, and membership of the value is closure of $\mathcal M_\sigma$ under the
operation $f,a\mapsto f(a)$, which holds because $f(a)\in\mathcal M_\sigma$ by definition of
$\mathcal M_{\tau'\to\sigma}$).

For (M$\le$) and (M$<$) we use that $p'=\interp{P}_{\rho'}\in\mathcal M_{\tau'\to\sigma}$
(IH (M0) for $P$) is non-strictly monotone in its argument:
\begin{equation}\label{eq:nmP}
  q\le_{\tau'}q'\ \Longrightarrow\ p'(q)\le_\sigma p'(q').
\end{equation}
This is now \emph{immediate from membership} ($p'\in\mathcal M_{\tau'\to\sigma}$ means
exactly that $p'$ is both strictly and non-strictly monotone), with no auxiliary term
required.
\begin{itemize}
\item (M$\le$). By IH \textup{(M$\le$)} for $P$, $p\le_{\tau'\to\sigma}p'$, so pointwise at
  $q$, $p(q)\le_\sigma p'(q)$. By IH \textup{(M$\le$)} for $Q$, $q\le_{\tau'}q'$, so by
  \eqref{eq:nmP}, $p'(q)\le_\sigma p'(q')$. Transitivity of $\le$ gives
  $p(q)\le_\sigma p'(q')$, i.e.\ $\interp{PQ}_\rho\le_\sigma\interp{PQ}_{\rho'}$.
\item (M$<$). Either $x\in\FV(Q)$ or $x\in\FV(P)$ (with $\rho=\rho'$ off $x$).
  \begin{itemize}
  \item If $x\in\FV(Q)$: by IH \textup{(M$<$)} for $Q$, $q<_{\tau'}q'$; since $p'$ is
    \emph{strictly} monotone (it lies in $\mathcal M_{\tau'\to\sigma}$), $p'(q)<_\sigma p'(q')$;
    and $p\le p'$ gives $p(q)\le_\sigma p'(q)$; mixed transitivity ($\le$ then $<$) gives
    $p(q)<_\sigma p'(q')$.
  \item If $x\notin\FV(Q)$ and $x\in\FV(P)$: then $q=q'$, and IH \textup{(M$<$)} for $P$
    gives $p<_{\tau'\to\sigma}p'$, so pointwise at $q$, $p(q)<_\sigma p'(q)=p'(q')$.
  \end{itemize}
  Either way $\interp{PQ}_\rho<_\sigma\interp{PQ}_{\rho'}$.
\end{itemize}
\end{proof}

\begin{remark}[The key simplification]\label{rem:keyfix}
Non-strict monotonicity of an interpreted functional in its argument is exactly membership
$\interp{M}_\rho\in\mathcal M_\sigma$, by the definition of the arrow domains
(Definition~\ref{def:dom}, Remark~\ref{rem:bundle}). It is \emph{not} the (false) statement
that every strictly monotone functional is non-strictly monotone (Remark~\ref{rem:leq}): we
do not claim that, and we do not need it. The earlier strategy of deriving non-strict
argument-monotonicity from environment-monotonicity at a fresh variable $w$ ran into a real
difficulty, because if one only demands strict monotonicity for membership, then a bound
variable used in function position could require the non-strict monotonicity of an
\emph{arbitrary} argument $a$, which need not hold. Bundling both monotonicities into
membership eliminates this difficulty at the source: every value (including a bound
argument) is non-strictly monotone by fiat, so the induction goes through cleanly and the
auxiliary-term detour is unnecessary.
\end{remark}

\section{Substitution, redex drop, congruence, and the main theorem}

\begin{lemma}[Substitution]\label{lem:subst}
Let $M:\sigma$, $x:\tau$, $N:\tau$, and let $\rho$ be an environment on $\FV(M[x:=N])$.
Then $\interp{M[x:=N]}_\rho=\interp{M}_{\rho[x:=\interp{N}_\rho]}$.
\end{lemma}

\begin{proof}
Induction on $M$, with $\alpha$-conventions so bound variables of $M$ avoid $x$ and
$\FV(N)$. We use that $\interp{T}_\rho$ depends only on $\rho{\restriction}\FV(T)$, and that
$\interp{N}_\rho\in\mathcal M_\tau$ (Lemma~\ref{lem:official}, (M0)), so
$\rho[x:=\interp N_\rho]$ is again an environment.
\begin{itemize}
\item $M=x$: both sides $=\interp N_\rho$. \quad $M=y\neq x$: both sides $=\rho(y)$.
\item $M=PQ$: by IH, $\interp{(PQ)[x:=N]}_\rho=\interp{P[x:=N]}_\rho(\interp{Q[x:=N]}_\rho)
  =\interp{P}_{\rho[x:=\interp N_\rho]}(\interp{Q}_{\rho[x:=\interp N_\rho]})
  =\interp{PQ}_{\rho[x:=\interp N_\rho]}$.
\item $M=\lambda z.P$, $z\notin\{x\}\cup\FV(N)$. For all $a$,
  \begin{align*}
   \interp{(\lambda z.P)[x:=N]}_\rho(a)
   &=\mathrm{add}\bigl(\interp{P[x:=N]}_{\rho[z:=a]},1{+}\|a\|\bigr)\\
   &\overset{\text{IH}}{=}\mathrm{add}\bigl(\interp{P}_{\rho[z:=a][x:=\interp N_{\rho[z:=a]}]},1{+}\|a\|\bigr)\\
   &=\mathrm{add}\bigl(\interp{P}_{\rho[x:=\interp N_\rho][z:=a]},1{+}\|a\|\bigr)
   =\interp{\lambda z.P}_{\rho[x:=\interp N_\rho]}(a),
  \end{align*}
  using $\interp N_{\rho[z:=a]}=\interp N_\rho$ ($z\notin\FV(N)$) and that updates at
  $z\neq x$ commute. Equality for all $a$ gives equality of the functions.
\end{itemize}
\end{proof}

\begin{lemma}[Strict drop at a redex]\label{lem:redex}
For every redex $(\lambda x^\tau.M)\,N$ of type $\sigma$ and every environment $\rho$ on its
free variables, $\interp{M[x:=N]}_\rho<_\sigma\interp{(\lambda x^\tau.M)\,N}_\rho$.
\end{lemma}

\begin{proof}
Let $a:=\interp N_\rho\in\mathcal M_\tau$; then $\|a\|_\tau\ge1$, so $1+\|a\|_\tau>0$. By
Definition~\ref{def:interp},
$\interp{(\lambda x.M)N}_\rho=\interp{\lambda x.M}_\rho(a)=
\mathrm{add}_\sigma(\interp{M}_{\rho[x:=a]},1{+}\|a\|_\tau)$, and by Lemma~\ref{lem:subst},
$\interp{M[x:=N]}_\rho=\interp{M}_{\rho[x:=a]}$. By Lemma~\ref{lem:add}(d),
\[
  \interp{M[x:=N]}_\rho=\mathrm{add}_\sigma(\interp{M}_{\rho[x:=a]},0)
  <_\sigma\mathrm{add}_\sigma(\interp{M}_{\rho[x:=a]},1{+}\|a\|_\tau)=\interp{(\lambda x.M)N}_\rho.\qedhere
\]
\end{proof}

\begin{remark}[Robustness]
The drop comes entirely from the strictly positive shift at the abstraction.
\emph{Discarding} ($x\notin\FV(M)$): $\interp{M}_{\rho[x:=a]}=\interp{M}_\rho$ is independent
of $N$, yet $\mathrm{add}(\interp M_\rho,0)<\mathrm{add}(\interp M_\rho,1{+}\|a\|)$.
\emph{Duplicating} ($x$ occurs $\ge2$ times): the Substitution Lemma uses the single value
$a=\interp N_\rho$ wherever $x$ occurs, so no copy of $N$ is re-evaluated and the measure
cannot blow up.
\end{remark}

\begin{lemma}[Congruence]\label{lem:cong}
Suppose $M,M':\sigma$ satisfy $\interp{M}_\rho<_\sigma\interp{M'}_\rho$ for every
environment $\rho$. Then for every such $\rho$:
\textup{(i)} $\interp{M\,P}_\rho<\interp{M'\,P}_\rho$;
\textup{(ii)} $\interp{P\,M}_\rho<\interp{P\,M'}_\rho$;
\textup{(iii)} $\interp{\lambda y.M}_\rho<\interp{\lambda y.M'}_\rho$.
\end{lemma}

\begin{proof}
\emph{(i)} $M,M':\tau'\to\sigma$, $P:\tau'$. The hypothesis is
$\interp{M}_\rho(b)<_\sigma\interp{M'}_\rho(b)$ for all $b$; take $b=\interp P_\rho$.

\emph{(ii)} $P:\tau'\to\sigma$, $M,M':\tau'$. By hypothesis
$\interp{M}_\rho<_{\tau'}\interp{M'}_\rho$. Since $\interp P_\rho\in\mathcal M_{\tau'\to\sigma}$
is strictly monotone (Lemma~\ref{lem:official}, (M0)),
$\interp{PM}_\rho=\interp P_\rho(\interp M_\rho)<_\sigma\interp P_\rho(\interp{M'}_\rho)=\interp{PM'}_\rho$.

\emph{(iii)} For $y:\tau'$ and every $a$, the hypothesis at $\rho[y:=a]$ gives
$\interp{M}_{\rho[y:=a]}<_\sigma\interp{M'}_{\rho[y:=a]}$, and Lemma~\ref{lem:add}(c) yields
$\interp{\lambda y.M}_\rho(a)<_\sigma\interp{\lambda y.M'}_\rho(a)$ for all $a$, hence
$\interp{\lambda y.M}_\rho<_{\tau'\to\sigma}\interp{\lambda y.M'}_\rho$.
\end{proof}

\begin{proposition}[Strict decrease under one step]\label{prop:step}
If $M\To N$ then $M,N:\sigma$ for a common $\sigma$, and for every environment
$\rho$ on $\FV(M)\supseteq\FV(N)$, $\interp{N}_\rho<_\sigma\interp{M}_\rho$.
\end{proposition}

\begin{proof}
Induction on the context of the contracted redex.
\emph{Head:} $M=(\lambda x.P)Q\To P[x:=Q]=N$ -- Lemma~\ref{lem:redex}.
\emph{Function position:} $M=M_1P\To M_1'P$ with $M_1\To M_1'$ -- IH gives
$\interp{M_1'}_\rho<\interp{M_1}_\rho$ for all $\rho$; apply Lemma~\ref{lem:cong}(i).
\emph{Argument position:} $M=PM_1\To PM_1'$ -- IH and Lemma~\ref{lem:cong}(ii).
\emph{Under $\lambda$:} $M=\lambda y.M_1\To\lambda y.M_1'$ -- IH and
Lemma~\ref{lem:cong}(iii). Since $\FV(N)\subseteq\FV(M)$, restricting $\rho$ to $\FV(N)$ is
legitimate.
\end{proof}

\section{The ordinal assignment}

\begin{definition}[The measure $\#$]\label{def:hash}
Let $\rho_0$ be the canonical environment, $\rho_0(x):=\mathbf 0_\tau$ for $x:\tau$. For
$M:\sigma$ set $\#M:=\|\interp{M}_{\rho_0}\|_\sigma\in\mathbb N_{\ge1}\subseteq\Ord$.
\end{definition}

This is well defined: each $\mathbf 0_\tau=\mathrm{lift}_\tau(1)\in\mathcal M_\tau$ exists
(Lemma~\ref{lem:basic}(a)), so $\rho_0$ is an environment; then
$\interp{M}_{\rho_0}\in\mathcal M_\sigma$ by Lemma~\ref{lem:official}, and the weight is a
positive integer. It is a finite structural recursion mentioning neither reduction nor
normalization.

\begin{theorem}[Main]\label{thm:main}
For all $M,N\in\Lambda_\to$, $M\To N\Rightarrow\#M>\#N$.
\end{theorem}

\begin{proof}
$M\To N$ gives $M,N:\sigma$ with $\FV(N)\subseteq\FV(M)$, so $\rho_0$ is defined on both.
By Proposition~\ref{prop:step}, $\interp{N}_{\rho_0}<_\sigma\interp{M}_{\rho_0}$; applying
the strictly order-preserving weight (Lemma~\ref{lem:basic}(b)),
$\#N=\|\interp{N}_{\rho_0}\|_\sigma<\|\interp{M}_{\rho_0}\|_\sigma=\#M$.
\end{proof}

\begin{corollary}[Strong normalization]\label{cor:SN}
The simply typed $\lambda$-calculus is strongly normalizing.
\end{corollary}

\begin{proof}
An infinite reduction $M_0\To M_1\To\cdots$ would, by Theorem~\ref{thm:main}, give an
infinite strictly decreasing sequence $\#M_0>\#M_1>\cdots$ of positive integers,
contradicting well-foundedness of $(\mathbb N,<)\subseteq(\Ord,<)$. Since $\#$ was built by
pure structural recursion (Definitions~\ref{def:dom}--\ref{def:hash}) and
Theorem~\ref{thm:main} was proved without assuming termination, this is a genuine,
non-circular, direct proof of SN via well-foundedness of $\Ord$.
\end{proof}

\section*{Discussion}

The map $\#$ is the strictly monotone hereditary functional interpretation of
Gandy/de Vrijer, recast so that every comparison is mediated by an explicit natural-number
weight $\|\cdot\|_\sigma$. Three design points make it work:
\begin{enumerate}
\item the shift $1+\|a\|_\sigma$ at each abstraction (Definition~\ref{def:interp}), which
  makes abstractions strictly monotone \emph{including the unused-variable case}
  (Remark~\ref{rem:why}) and supplies the strict drop at a redex (Lemma~\ref{lem:redex});
\item restricting each function space to functionals that are \emph{both strictly and
  non-strictly monotone} (Definition~\ref{def:dom}, Remark~\ref{rem:bundle}), a class that
  is preserved by $\mathrm{lift}$, $\mathrm{add}$ and the interpretation
  (Lemmas~\ref{lem:basic},~\ref{lem:add},~\ref{lem:official}); strict monotonicity
  propagates the strict drop through reduction in argument position
  (Lemma~\ref{lem:cong}(ii)), while non-strict monotonicity --- now bundled into
  membership, hence available for \emph{every} value including bound arguments --- is what
  makes the environment/argument induction go through (cf.\ Remark~\ref{rem:keyfix}); and
\item the \emph{pointwise} non-strict order $\le_\sigma$ (Definition~\ref{def:dom},
  Remark~\ref{rem:leq}) as the inductive invariant. The naive reading ``$\le$ means $<$ or
  $=$'' is \emph{false} at arrow types and would break the induction.
\end{enumerate}
An earlier attempt used a fixed ``$+1$'' shift; that interpretation is not well defined,
because $\lambda x.M$ with $x\notin\FV(M)$ would be constant, breaking
Lemma~\ref{lem:cong}(ii). A second subtlety, repaired here, is that requiring only
\emph{strict} monotonicity for membership leaves a gap: non-strict monotonicity of an
interpreted functional was then needed for an \emph{arbitrary} bound argument, which a
merely strictly monotone functional need not satisfy. Putting both monotonicities into the
definition of $\mathcal M_{\sigma\to\tau}$ closes this gap directly, yielding a complete,
non-circular proof. All critical inequalities have additionally been verified numerically on
many multi-step reductions of higher order, with no violation.
\end{document}
